\chapter{Resultados y comparación}
\label{ch:ch7_results}

Este capítulo reúne lo que el trabajo establece y lo separa, sin
suavizarlo, de lo que no establece. El motor se ha modelado, escrito,
verificado en simulación frente a vectores oficiales y sintetizado sobre
una biblioteca de celdas estándar real; todavía no ha funcionado en la
\ac{FPGA}. Por eso cada cifra que sigue lleva su marca de procedencia, y
la campaña de entropía --- el único resultado que necesita silicio en el
banco --- se presenta como un protocolo por ejecutar, no como datos.

\section{Estado frente a los objetivos}
\label{sec:status}

La Tabla~\ref{tab:status} relaciona los objetivos de la
Sección~\ref{sec:objectives} con las evidencias.

\begin{table}[htbp]
\centering
\caption{Objetivos y evidencia que respalda cada uno.}
\label{tab:status}
\small
\begin{tabular}{@{}llp{7.4cm}@{}}
\toprule
Obj. & Estado & Evidencia \\
\midrule
O1 & cumplido  & Anillo a nivel de transistor en el \ac{PDK} abierto IHP SG13G2,
              con el ruido de sus propios modelos de dispositivo
              (Sección~\ref{sec:ihp_ring}); pendiente el \ac{PDK} de fundición \\
O2 & cumplido  & Jitter extraído de las formas de onda; térmico y $1/f$
              separados (banco y función de sensibilidad al impulso);
              estimador validado sobre datos sintéticos con $Q$ conocido \\
O3 & cumplido  & Modelo estocástico cerrado: cálculo exacto con la gaussiana envuelta,
              punto de diseño $Q \ge 0{,}2286$ (Sección~\ref{sec:opoint}) \\
O4 & cumplido  & Motor escrito en VHDL-2008 portable; siete bancos de pruebas pasan
              con los vectores de FIPS-197, SP 800-38B, RFC 4493 y \ac{CAVP} \\
O5 & parcial   & El firmware del verificador compila y el protocolo está cerrado;
              falta la ejecución en placa \\
O6 & pendiente & Requiere la \ac{FPGA} en el banco
              (Sección~\ref{sec:protocol}) \\
O7 & cumplido  & Síntesis en IHP SG13G2, área por bloque y vía de optimización
              (Capítulo~\ref{ch:ch6_asic}) \\
O8 & parcial   & Comparación estructural cerrada (Sección~\ref{sec:se_comparison});
              el lado ESP32 de la comparación numérica está medido según
              SP 800-90B (Sección~\ref{sec:esp32_90b}); el lado del motor necesita O6 \\
\bottomrule
\end{tabular}
\end{table}

El resultado mínimo defendible declarado en el Capítulo~\ref{ch:ch1_intro} ---
O1--O3 junto con O5 --- se alcanza en su parte de modelado y diseño.
Lo que aún le falta a O5 es la puesta en marcha, no el diseño.

\section{Ruido del anillo a nivel de transistor}
\label{sec:spice_results}

El anillo se simuló con una fuente de ruido de corriente de densidad
espectral controlada inyectada en el nodo de conmutación, barriendo la
densidad y midiendo la desviación típica de los cruces de flanco de
subida a lo largo de 500 periodos. La Tabla~\ref{tab:spice} da el
resultado para un anillo de cinco etapas con una carga de 10~fF.

\begin{table}[htbp]
\centering
\caption{Jitter por periodo del anillo de cinco etapas frente a la densidad
de ruido inyectada \simu{}. $\sigma_{\mathrm{rel}} = \sigma/T$.}
\label{tab:spice}
\small
\begin{tabular}{@{}lrrrr@{}}
\toprule
Densidad de ruido (A$^2$/Hz) & Cruces & $T$ (ps) & $\sigma$ (fs) &
$\sigma_{\mathrm{rel}}$ \\
\midrule
ninguno (suelo numérico) & 500 & 230,30 & 1,38   & $6{,}0\times10^{-6}$ \\
$1\times10^{-20}$        & 501 & 230,37 & 1200   & $5{,}2\times10^{-3}$ \\
$4\times10^{-20}$        & 507 & 229,98 & 2416   & $1{,}1\times10^{-2}$ \\
$1{,}6\times10^{-19}$    & 509 & 227,77 & 25588  & $1{,}1\times10^{-1}$ \\
\bottomrule
\end{tabular}
\end{table}

Tres lecturas, y la tercera es la que importa.

\paragraph{El suelo numérico queda lo bastante por debajo de la señal.} Con
la fuente de ruido apagada, la medida devuelve 1,38~fs, que es la
resolución de la propia simulación, no física. Está unas 870 veces por
debajo del menor jitter que hay que medir, así que el suelo no contamina
el barrido. Llegar a él exigió cuidado: el primer intento devolvió
0,33~ps, del mismo orden que la magnitud estudiada, y la causa era que el
transitorio de arranque del anillo entraba en la estadística. Descartar
los veinte primeros ciclos lo eliminó. Un suelo de medida que coincide con
la señal no se distingue de un resultado, y precisamente por eso el suelo
se caracterizó antes del barrido y no después.

\paragraph{La ley de escala se cumple donde debe.} El jitter térmico se
acumula como la raíz cuadrada de la densidad de ruido, así que cuadruplicar
la densidad debe duplicar $\sigma$. De $1\times10^{-20}$ a
$4\times10^{-20}$ la medida da un factor de 2,01 frente a los 2,00
previstos. Esa coincidencia es la prueba de que lo que se mide es el ruido
inyectado y no un artefacto del resolvedor.

\paragraph{Y se rompe, a la vista, en lo alto del barrido.} Con
$1{,}6\times10^{-19}$, dieciséis veces la densidad de referencia, $\sigma$
crece en un factor de 21,3 donde la ley predice 4,0, y el periodo medio
baja de 230,4 a 227,8~ps. Una fuente de ruido de media nula no puede
desplazar un periodo medio. Por tanto, el punto queda fuera del régimen
que describe el modelo --- la perturbación es lo bastante grande como para
cambiar el punto de funcionamiento de los inversores --- y se presenta en
lugar de descartarse, porque el límite de validez de un modelo es en sí
mismo un resultado. Solo los dos puntos centrales se usan en lo que sigue.

\section{El anillo con los transistores de un PDK abierto}
\label{sec:ihp_ring}

El mismo anillo de cinco etapas se simuló con los transistores de baja
tensión del \ac{PDK} abierto IHP SG13G2 ($W_n = 0{,}5~\mu$m,
$W_p = 1{,}2~\mu$m, $L = 0{,}13~\mu$m, 1,2~V, 27~$^\circ$C, esquina
típica, 1~fF por nodo como estimación del cableado previa al layout), con
la densidad de ruido medida en los mismos dispositivos
(Sección~\ref{sec:spice}) inyectada en cada nodo. La
Tabla~\ref{tab:ihp_ring} da el resultado sobre 400~ns, unos 1070
periodos por punto.

\begin{table}[htbp]
\centering
\caption{Anillo de cinco etapas en IHP SG13G2 \simu{}. La densidad de ruido
no se supone: es el suelo blanco del propio modelo de ruido de los
transistores en el punto de conmutación. La potencia es la media sobre
periodos completos.}
\label{tab:ihp_ring}
\small
\begin{tabular}{@{}lrrrrr@{}}
\toprule
Ruido inyectado & $T$ (ps) & $\sigma$ (fs) & $\sigma/T$ & Exponente & $P$ ($\mu$W) \\
\midrule
ninguno (suelo numérico)        & 366,04 & 0,70   & $1{,}9\times10^{-6}$ & 0,47 & 143,0 \\
térmico del \ac{PDK}, $4{,}13\times10^{-24}$~A$^2$/Hz & 365,97 & 180,37 & $4{,}93\times10^{-4}$ & 1,05 & 143,0 \\
térmico $\times4$               & 365,97 & 358,64 & $9{,}80\times10^{-4}$ & 1,08 & 143,0 \\
térmico + $1/f$ del \ac{PDK}    & 365,97 & 183,33 & $5{,}01\times10^{-4}$ & 1,01 & 143,0 \\
\bottomrule
\end{tabular}
\end{table}

Primero, la resolución: 0,70~fs, 259 veces por debajo del jitter que
produce el ruido de los dispositivos, y con un exponente de 0,47 que la
identifica como ruido numérico y no como un paseo aleatorio. Cuadruplicar
la densidad duplica el jitter con un 0,6\,\% de diferencia
($\times1{,}988$), y el periodo medio no se mueve, así que el banco se
mantiene en el régimen lineal que el banco provisional abandonaba en su
punto más alto. El exponente de acumulación es 1,05: ruido blanco, como
supone el modelo estocástico. Como comprobación independiente, la
expresión cerrada de Abidi para osciladores en
anillo~\cite{abidi_2006_ring}, evaluada con magnitudes medidas en el
mismo punto de polarización, predice $\sigma/T = 1{,}2\times10^{-3}$; la
simulación da 0,42 veces ese valor, el mismo orden. La forma de canal
largo de esa fórmula, que pasa por la tensión umbral, no se aplica aquí:
en el punto de conmutación los dispositivos están solo 0,07~V por encima
del umbral.

Con $\sigma/T = 4{,}93\times10^{-4}$, el divisor para $Q \ge 0{,}2286$ es
$K_D = 941\,107$, un intervalo de muestreo de 344~$\mu$s con un anillo de
muestreo idéntico, y 2,9~kbit/s por \ac{ERO}; si además se cuenta el
jitter independiente del anillo de muestreo, $K_D$ se reduce a la mitad y
la tasa se duplica, pero eso exige que los dos anillos no compartan ruido
a través de la alimentación o del sustrato, algo que solo el silicio
puede mostrar. Una consecuencia de diseño se aplicó de inmediato: el
contador del divisor tenía 20 bits, así que el mayor $K_D$ era
$2^{20} = 1\,048\,576$, un margen del 11\,\% sobre lo requerido en la
esquina típica a temperatura ambiente. Ahora tiene 24 bits (cuatro
flip-flops, 0,04~kGE), y el banco del motor completo y su verificación
cruzada se volvieron a ejecutar con él.

\paragraph{Ruido de parpadeo.} Añadir el ruido $1/f$ del \ac{PDK} no
cambia nada en 400~ns: la separación $E[D^2] = a_1 M + a_2 M^2$, validada
antes sobre fases sintéticas (componente térmica recuperada al 0,1\,\%
con y sin $1/f$), da un $a_2$ compatible con cero. Amplificar la fuente
$1/f$ para hacerla visible no funciona, y el motivo es instructivo: la
fuente es de banda ancha y, amplificada mil veces, su densidad cerca de la
frecuencia de oscilación ya es siete veces el suelo térmico, así que el
anillo la convierte como ruido blanco y el término lineal crece en el
factor esperado mientras que el término cuadrático no escala. La
respuesta llega, en cambio, de la función de sensibilidad al impulso de
Hajimiri y Lee~\cite{hajimiri_1998_general}, medida inyectando el
1\,\% de la carga del nodo en 40 fases (lineal al 0,1\,\% entre 0,5 y
2\,\%). Es casi simétrica, $|\Gamma_{dc}/\Gamma_{rms}| \approx 0{,}03$--$0{,}05$,
así que el ruido de baja frecuencia inyectado de forma uniforme a lo largo
del periodo apenas se convierte. Pero los transistores no son simétricos:
medidos por separado, el dispositivo de canal p tiene tres veces el ruido
$1/f$ del de canal n, y cada uno inyecta ruido sobre todo mientras
conduce, durante uno de los dos lóbulos opuestos. Los dos límites acotan
la realidad: con ruido estacionario, como en el banco, la frecuencia de
esquina $1/f^3$ es 36~kHz y el parpadeo iguala al jitter térmico tras
27\,000 periodos; con cada dispositivo confinado a su lóbulo, la
frecuencia de esquina es 5,3~MHz y el cruce, 185 periodos. En ambos
casos, con el divisor de funcionamiento el parpadeo domina la varianza
acumulada, por un factor de entre 35 y 5\,000. De ahí se siguen dos
consecuencias. El jitter que fija la cota de entropía debe medirse con
acumulación corta, por debajo de unos cien periodos para este anillo, y
separarse del término cuadrático --- deducirlo del jitter acumulado total
en $K_D$ atribuiría una entropía que no existe. Y el banco estacionario
puede sobrestimar la varianza térmica hasta en un factor dos, lo que
duplicaría $K_D$; el contador de 24 bits lo cubre, y la herramienta que lo
resuelve es un análisis de ruido transitorio cicloestacionario, un uso
concreto de un simulador comercial.

\section{Punto de funcionamiento: divisor y tasa de bits}
\label{sec:opoint}

El modelo estocástico fija el objetivo. Para los umbrales PTG.2 de AIS
20/31 v3.0, el criterio de min-entropía ($H_\infty \ge 0{,}98$) exige
$Q \ge 0{,}2286$ y es más exigente que el criterio de Shannon
($H \ge 0{,}9998$, $Q \ge 0{,}2022$); por eso el diseño apunta al primero,
con $Q = 0{,}30$ como valor de consigna de trabajo, que da
$H_\infty = 0{,}9951$. Dos propiedades del modelo se demostraron en lugar
de suponerse: la serie truncada de la ecuación~14
de~\cite{baudet_2011_security} \emph{sobrestima} la entropía y, por
tanto, no es una cota inferior, así que se usa en su lugar el cálculo
exacto con la gaussiana envuelta; y la deducción de la min-entropía con el
primer armónico coincide con la exacta en seis cifras en todo el rango de
interés.

Como $Q_{\mathrm{gen}} = K_D \cdot Q_{\mathrm{bruto}}$, el divisor se
deduce del jitter que mida la campaña. La Tabla~\ref{tab:opoint} da los
puntos de funcionamiento para el rango de jitter esperable en esta
tecnología~\cite{petura_2016_survey}, suponiendo RO2 a 500~MHz.

\begin{table}[htbp]
\centering
\caption{Divisor y tasa de bits resultante por anillo para
$Q \ge 0{,}2286$ \est{}.}
\label{tab:opoint}
\small
\begin{tabular}{@{}rrrr@{}}
\toprule
$\sigma/T$ por periodo & $Q_{\mathrm{bruto}}$ & $K_D$ & Tasa de bits \\
\midrule
$7{,}1\times10^{-4}$ & $5\times10^{-7}$ & 457\,200 & 1,09~kbit/s \\
$1{,}4\times10^{-3}$ & $2\times10^{-6}$ & 114\,300 & 4,37~kbit/s \\
$2{,}8\times10^{-3}$ & $8\times10^{-6}$ & 28\,575  & 17,5~kbit/s \\
\bottomrule
\end{tabular}
\end{table}

Estos valores encierran el punto de diseño publicado por Fischer y
Lubicz~\cite{fischer_2014_embedded} ($K_D \approx 430\,000$ para
$\sigma = 5{,}01$~ps sobre $T = 7{,}81$~ns), y esa es la comprobación de
coherencia de que el modelo se aplica como pretendían sus autores.

La consecuencia para la arquitectura se dice sin rodeos: un solo \ac{ERO}
entrega kbit/s, no Mbit/s. Ese es el precio de una entropía demostrable,
y aquí se puede pagar porque el producto del motor es una semilla de
256 bits para el \ac{DRBG}, no un flujo de bits. Más tasa de bits
significa más anillos independientes, cada uno con su propia cota, no un
anillo más rápido.

\section{Verificación funcional consolidada}
\label{sec:verif_results}

El Capítulo~\ref{ch:ch5_verification} describe la estrategia; esta
sección recoge los totales. Cada cifra de la Tabla~\ref{tab:verif} se
obtuvo ejecutando la batería de pruebas, no leyendo un registro anterior.

\begin{table}[htbp]
\centering
\caption{Totales de la verificación \simu{}.}
\label{tab:verif}
\small
\begin{tabular}{@{}lp{7.6cm}l@{}}
\toprule
Batería & Qué comprueba & Resultado \\
\midrule
Bancos de pruebas VHDL (GHDL) & anillos, \ac{ERO}, captura, tests de salud,
  \ac{I2C}, \ac{AES}, \ac{CMAC}, \ac{CTRDRBG}, nivel superior & 7/7 pasan \\
\texttt{ctr\_drbg\_ref.py} & vectores \ac{CAVP}, con y sin
  función de derivación & 960/960 \\
\texttt{ero\_model.py} & modelo estocástico, exactitud y cotas & 24/24 \\
\texttt{jitter\_estimator.py} & recuperación del jitter sobre datos
  sintéticos de $Q$ conocido & 20/20 \\
\texttt{cortes\_salud.py} & umbrales de los tests de salud frente a las
  fórmulas de SP 800-90B & 6/6 \\
\texttt{verifica\_top.py} & verificación cruzada del banco de nivel
  superior con un \ac{AES} independiente & 5/5 \\
\bottomrule
\end{tabular}
\end{table}

Dos de ellas merecen una frase. Los 960 casos \ac{CAVP} son la prueba
individual más sólida del trabajo: el generador no solo es coherente
consigo mismo, sino que reproduce las salidas que el \ac{NIST} publica
para \ac{CTRDRBG} con AES-128, tanto en la configuración con función de
derivación como en la de entropía en bruto. Y los umbrales de los tests de
salud, $C_{RCT} = 22$ y $C_{APT} = 596$, se habían calculado a mano a
partir de la norma y nunca se habían comprobado; volver a obtenerlos con
las fórmulas de SP 800-90B~\cite{turan_2018_sp80090b} confirma ambos. La
misma comprobación cuantifica el coste de la prueba del complemento
añadida al \ac{APT}, que es una aportación de este trabajo y no de la
norma: eleva la tasa de falsas alarmas por ventana de $2^{-20{,}15}$ a
$2^{-20{,}14}$, un aumento del 0,8\,\%, porque la cola complementaria está
a 5,7 desviaciones típicas. Por tanto, el umbral sigue siendo válido tal
como está y la adición sale gratis en la práctica.

\section{Comparación con elementos seguros comerciales}
\label{sec:se_comparison}

El motor ocupa el mismo hueco funcional que un elemento seguro comercial:
un pequeño dispositivo acompañante que genera y guarda claves para un
microcontrolador anfitrión a través de un bus de dos hilos. La
Tabla~\ref{tab:se} lo sitúa junto a los tres componentes que dominan el
mercado de elementos seguros \ac{I2C}.

\begin{table}[htbp]
\centering
\caption{El motor frente a elementos seguros comerciales. Datos de
certificación tomados de las hojas de datos de los fabricantes y de los
registros de validación del NIST~\cite{microchip_atecc608b_ds, nist_esv_e46, nxp_se050_ds,
infineon_optiga_trustm_ds}.}
\label{tab:se}
\small
\begin{tabular}{@{}lp{2.6cm}p{2.4cm}p{2.2cm}p{2.6cm}@{}}
\toprule
 & ATECC608B & SE050 & OPTIGA Trust M & Este trabajo \\
\midrule
Bus            & \ac{I2C} & \ac{I2C} & \ac{I2C} & \ac{I2C} \\
Certificación  & JIL High; NIST ESV \#E46 para la fuente de entropía & CC EAL6+ (hasta el SO) & CC EAL6+ (hardware) & ninguna \\
Fuente de entropía & osciladores en anillo, según el registro ESV & ``TRNG, SP~800-90B'' en la hoja de datos & no indicada & \ac{ERO}, publicada \\
Entropía declarada & 0,5071~bit/muestra, validada por el NIST & no publicada & no publicada & cota obtenida del jitter medido \\
Modelo estocástico & no publicado & no publicado & no publicado & publicado \\
Exportación de claves & no es posible & no es posible & no es posible & solo en modo de prueba \\
Inspeccionable & no & no & no & sí, \ac{RTL} publicado \\
\bottomrule
\end{tabular}
\end{table}

La lectura honesta de esa tabla tiene dos mitades.

En prestaciones y nivel de garantía ganan los componentes comerciales, y
con mucha diferencia. Dos de ellos tienen evaluaciones CC EAL6+, que
suponen un grado de escrutinio por terceros --- incluida la resistencia a
ataques físicos --- que ningún proyecto académico reproduce. Ofrecen
criptografía asimétrica, flujos de aprovisionamiento certificados y
soporte para la fabricación. Este trabajo ofrece generación de claves
simétricas y autenticación, y no tiene ninguna certificación.

En lo que de verdad trata este TFM, la comparación tiene más matices que
``abierto frente a cerrado'', y conviene ser exacto. De los tres, solo el
ATECC608B tiene su fuente de entropía validada según SP~800-90B: el
registro del NIST cita osciladores en anillo y declara 0,5071~bit de
entropía por muestra~\cite{nist_esv_e46}. Es una cifra real, de un
tercero, y hay que reconocerla. Lo que ningún fabricante publica ---
tampoco Microchip --- es el modelo estocástico que hay detrás de la cifra,
la medida de jitter que la respalda ni el propio circuito; la hoja de
datos de un competidor afirma cumplir la norma sin un registro de
validación y la del otro no menciona SP~800-90 en absoluto. Se confía en
su aleatoriedad porque se confía en el fabricante y en el esquema de
evaluación. El motor descrito aquí hace el intercambio contrario: no tiene
certificado, y puede mostrar de dónde viene cada bit de entropía, qué
medida física lo acota y en qué se convierte la cota cuando cambian la
temperatura o el divisor. Es la misma distinción con la que tropieza,
desde el otro lado, el estudio de caja negra del generador del ESP32 de
este trabajo (Sección~\ref{sec:esp32_90b}) --- una caja negra se puede
probar estadísticamente,
pero no se puede acotar --- y es la razón por la que este motor se
construyó como una caja blanca.

Cuál de los dos es preferible depende de la pregunta que se haga. Para
sacar un producto al mercado, el componente certificado. Para saber cuánto
vale la aleatoriedad de ese producto, solo responde el abierto.

